% ============================================================
%  FA-KGD -> MIDL 2027: ranked, literature-grounded revision plan
%  Internal planning document (not a submission).
%  Build:  pdflatex revision_plan && bibtex revision_plan &&
%          pdflatex revision_plan && pdflatex revision_plan
% ============================================================
\documentclass[10pt,a4paper]{article}
\usepackage[margin=2.2cm]{geometry}
\usepackage[T1]{fontenc}
\usepackage{lmodern}
\usepackage{microtype}
\usepackage{amsmath,amssymb}
\usepackage{booktabs}
\usepackage{tabularx}
\usepackage{array}
\usepackage{enumitem}
\usepackage[table]{xcolor}
\usepackage[numbers,sort&compress]{natbib}
\usepackage[colorlinks,linkcolor=blue!50!black,citecolor=green!35!black,urlcolor=blue!50!black]{hyperref}

\setlist{nosep,leftmargin=1.4em}
\newcolumntype{L}{>{\raggedright\arraybackslash}X}
\renewcommand{\arraystretch}{1.15}

\newcommand{\fakgd}{FA-KGD}
\newcommand{\pigdm}{$\Pi$GDM}
\pdfstringdefDisableCommands{\def\pigdm{PiGDM}\def\fakgd{FA-KGD}}
\newcommand{\code}[1]{\texttt{\small #1}}
\newcommand{\impact}[1]{\textbf{Impact:} #1}
\newcommand{\tier}[1]{\colorbox{gray!15}{\textsc{#1}}}

% One block per proposed change
\newcommand{\change}[2]{\subsection*{#1\quad #2}\addcontentsline{toc}{subsection}{\texorpdfstring{#1\quad #2}{#1 #2}}}
\newcommand{\field}[1]{\par\smallskip\noindent\textbf{#1}\enspace}

\title{\fakgd{} $\rightarrow$ MIDL 2027\\[2pt]
\large A ranked, literature-grounded revision plan}
\author{Internal planning document for the \fakgd{} author team}
\date{Prepared 5 October 2026\quad$\cdot$\quad Full-paper deadline 4 December 2026 (AoE)}

\begin{document}
\maketitle

% ------------------------------------------------------------
\section*{Executive summary}
% ------------------------------------------------------------
The CS231N report (\code{paper/cvpr2017AuthorKit/latex/egpaper\_final.tex}) contains one
result that is potentially interesting for MIDL. A diffusion sampler with isotropic
data consistency (DC) \emph{loses} PSNR as the number of reverse steps $T$ grows,
while a frequency-aware DC does not. As currently written, however, the paper would
likely be rejected. There are four reasons.
\begin{enumerate}
  \item Its physical premise, that k-space \emph{measurement} noise varies by 1--2 orders
        of magnitude across frequency, is incorrect for thermal noise
        \cite{roemer1990phased,kellman2005snr}. All headline numbers use synthetic
        radial noise.
  \item The baseline labelled ``\pigdm'' is not \pigdm~\cite{song2023pigdm}
        (\S\ref{sec:audit}).
  \item One of the three method components (the online M-step) is very likely inert
        in the configuration used.
  \item Closely related work from 2024--2026 is not discussed
        \cite{peng2024optimal,tian2026frequency,ozturkler2023smrd,lee2026msc}.
\end{enumerate}

None of these is fatal. Three things make the plan feasible: the deadline (8.5 weeks
away), one dedicated A100, and the existing multi-coil SENSE code path. The plan
reframes the contribution as a \textbf{frequency-resolved posterior-uncertainty view of
data consistency}. Under that view, \fakgd's ``per-frequency noise'' is reinterpreted as
a per-frequency \emph{prior/posterior variance}, and the claims become testable on real
multi-coil data. A hard go/no-go gate falls on \textbf{12 October 2026}
(\S\ref{sec:timeline}).

\begin{center}
\begin{tabularx}{\linewidth}{@{}c l L c@{}}
\toprule
Rank & ID & Change (highest $\rightarrow$ lowest impact) & Gate\\
\midrule
1 & C1 & Replace the noise premise with a posterior-variance interpretation; evaluate on real (prewhitened) noise only & G1\\
2 & C2 & Implement real \pigdm; rename the current baseline; tune every baseline on a calibration split & G1\\
3 & C3 & Diagnose the inert M-step; fix it or remove it & G1\\
4 & C4 & Make step-count degradation the central, mechanistically measured finding & G2\\
5 & C5 & Multi-coil SENSE with ESPIRiT maps and prewhitening as the primary operator & G2\\
6 & C6 & Plug-in evaluation across samplers, explicitly positioned against 2024--26 work & G2\\
7 & C7 & Leak-free splits, volume-level statistics, compute reporting & ---\\
8 & C8 & Anatomy-matched knee prior & ---\\
9 & C9 & Real low-SNR evidence (M4Raw 0.3\,T; fastMRI Prostate averages) & G3\\
10 & C10 & Pathology-region and hallucination analysis (fastMRI+) & ---\\
11 & C11 & Corrupted-sample robustness (spikes, motion lines) [stretch] & G3\\
12 & C12 & Downstream task evaluation (SKM-TEA / K2S) [post-deadline] & ---\\
13 & C13 & Posterior calibration of uncertainty [stretch] & ---\\
14 & C14 & Rewrite for MIDL: template, scope of claims, LLM policy & ---\\
\bottomrule
\end{tabularx}
\end{center}

\tableofcontents

% ------------------------------------------------------------
\section{Venue constraints (verified on 2027.midl.io, 5 Oct 2026)}
\label{sec:venue}
% ------------------------------------------------------------
\begin{center}
\begin{tabularx}{\linewidth}{@{}l L@{}}
\toprule
Item & Fact \cite{midl2027cfp}\\
\midrule
Dates & Abstract registration 30 Nov 2026; \textbf{full paper 4 Dec 2026}; reviews 10 Jan; rebuttal 18--24 Jan; notification 12 Feb 2027; camera-ready 12 Mar. Conference: Porto, 14--16 Jul 2027.\\
Main track & \textbf{10 pages} at submission (12 after rebuttal), excluding references, acknowledgements and appendices. ``Authors should justify the clinical or biomedical relevance of the problem.'' Reviewers are not required to read appendices.\\
Validation track & 14 pages; novelty not required (external validation, reader studies). This is not our track unless C1--C6 fail.\\
Review & \textbf{Single-blind} (since 2025) on OpenReview. Do not anonymise. Proceedings: PMLR.\\
Template & Official MIDL \LaTeX{} template (\url{github.com/MIDL-Conference/MIDLLatexTemplate}). Font, spacing and margins must not be changed.\\
Short papers & 3 pages, \emph{non-archival}. The track opens only after the full-paper cycle (about spring 2027), so it is a fallback, not a parallel option.\\
Prior posting & An arXiv preprint or non-archival workshop does not count as dual publication. The CS231N report is non-archival; to be safe, confirm with the PCs by email.\\
\textbf{LLM policy} & LLMs are allowed \textbf{only for minor language edits}. Generating ideas, analyses, results or substantive text is prohibited. \emph{Consequence:} the submission text must be written by the authors. This document is an internal planning aid; none of its prose should be pasted into the paper.\\
Acceptance & 61.8\% in 2023 (112/181, PMLR v227 preface); about 54\% in 2024. No official rate has been published for 2025--26.\\
\bottomrule
\end{tabularx}
\end{center}

\paragraph{What gets accepted.} Recent MIDL MRI-reconstruction papers
\cite{akoda2026consistency,deng2026subspace,karkalousos2023mtl} evaluate on fastMRI and/or
CMRxRecon across several acceleration factors. They compare against both diffusion and
unrolled baselines. They sell either speed (NFE) or physical consistency, and some add a
downstream task. This sets the minimum bar for C5--C7.

% ------------------------------------------------------------
\section{Audit of the current draft and code}
\label{sec:audit}
% ------------------------------------------------------------
Line numbers refer to \code{egpaper\_final.tex}; code paths are relative to the
repository root.

\begin{center}
\small
\begin{tabularx}{\linewidth}{@{}c L L@{}}
\toprule
Sev. & Finding & Evidence\\
\midrule
\cellcolor{red!15}A & The premise that k-space \emph{measurement} noise power varies 1--2 orders of magnitude between centre and periphery is physically wrong. Thermal noise is white per receive channel and correlated across coils \cite{roemer1990phased,pruessmann1999sense}. Receiver filtering changes noise along \emph{image position} on the readout axis and through the noise-equivalent bandwidth \cite{kellman2005snr}; it does not change per-sample k-space variance with spatial frequency. What does vary strongly with frequency is \emph{signal} power. & l.124--127; all sweeps use synthetic $\sigma^2(r)=\sigma_0^2(1+5(r/r_{\max})^2)$ (l.617--619, 755--757).\\
\cellcolor{red!15}A & The ``\pigdm'' baseline is a scalar-gain Kalman correction of the Tweedie estimate followed by fresh-noise re-noising. It has no vector--Jacobian product guidance on $x_t$ and no $r_t$ schedule, so it is not the method of \cite{song2023pigdm}. Eq.~(1) of the draft describes a different update from the one evaluated. & \code{src/samplers/pigdm.py:49--64}\\
\cellcolor{red!15}A & The M-step is very likely inert. $\alpha_t=1-(1-\alpha_0)\sigma_t/\sigma_{\max}\to1$ as $\sigma_t\to0$, so late steps barely update. In clamp mode, $\hat\sigma_i^2\le\hat\sigma^2_{i,\mathrm{init}}$, and early residuals exceed the initial value, so they are clamped back. This is consistent with the $\alpha$ sweep being identical to $10^{-3}$\,dB across all $\alpha_0$. & \code{src/samplers/fakgd.py:140--157}; Table~\code{tab:hparam} (l.1102--1104)\\
\cellcolor{orange!20}B & Hyperparameters are selected on the fastMRI validation split and results are reported on the same split. The test split is described as ``unused''. & l.594--596, 889--890\\
\cellcolor{orange!20}B & Overclaim: ``FA-KGD improves monotonically'' with $T$. Its absolute PSNR peaks at $T{=}50$ (30.24$\to$29.89\,dB; knee 31.92$\to$30.32). Only the \emph{gap} to the baseline grows. ``Scaling law'' (abstract) is not supported by four points. & l.81, 761--763, Table~\code{tab:tscaling}\\
\cellcolor{orange!20}B & Headline significance $p<10^{-30}$ treats 192 correlated slices as independent. The volume-level test ($n{=}12$, $p{=}2^{-12}$, i.e.\ every volume positive) is the defensible one. & l.777--779\\
\cellcolor{orange!20}B & The oracle ablation is internally inconsistent: $\eta=0.1$ vs.\ ``$\eta=0$''; ``three-slice'' vs.\ ``five-slice'' development set. Its PSNR values of about 60\,dB are far from the operating point. & l.908, 1132, 1150--1156\\
\cellcolor{orange!20}B & The paper describes three different noise sources: ACS coil-difference variance (Eq.~2), image background patches (multicoil), and synthetic radial noise. Readers cannot tell which drives which result. & l.384--399\\
\cellcolor{yellow!25}C & ``DDIM step'' in Alg.~1 and Fig.~1, but the code re-noises with fresh Gaussian noise (stochastic, not deterministic DDIM). & l.567; \code{fakgd.py:165--169}\\
\cellcolor{yellow!25}C & The knee results use the brain-trained prior. The \code{checkpoints/unet\_knee} files are supervised U-Nets, not diffusion priors. & l.640--641; \code{checkpoints/}\\
\cellcolor{yellow!25}C & The main experiments use a coil-combined single-channel operator, although a SENSE path exists in the code. & l.262--270; \code{src/samplers/multicoil.py}, \code{sense.py}\\
\bottomrule
\end{tabularx}
\end{center}

% ------------------------------------------------------------
\section{Position in the current literature}
\label{sec:lit}
% ------------------------------------------------------------
\begin{center}
\small
\begin{tabularx}{\linewidth}{@{}l L c c c@{}}
\toprule
Work & Where non-isotropy / adaptivity enters & Basis & Online & MRI\\
\midrule
\pigdm{} \cite{song2023pigdm}, DPS \cite{chung2023dps} & None: isotropic $r_t^2$ or Laplace approximation & --- & --- & no\\
Peng et al.\ \cite{peng2024optimal} & Posterior covariance of $x_0|x_t$, diagonal in an orthonormal basis; learned or offline Monte Carlo & DWT & no & no\\
TMPD \cite{boys2024tmpd}, MMPS \cite{rozet2024mmps} & Full Tweedie covariance via denoiser Jacobian (moment projection / CG with VJPs) & dense & yes & MMPS: knee\\
CA-DPS \cite{hamidi2024cadps} & Online diagonal covariance from score finite differences & pixel & yes & no\\
Schioppa \& Salimans \cite{schioppa2026covariance} & Fourier-structured Tweedie covariance (unconditional sampling) & Fourier & yes & no\\
Tian et al.\ \cite{tian2026frequency} & \textbf{Noise-conditioned frequency exposure} of measurements, returning to full-band DC & Haar/freq. & schedule & no\\
SMRD \cite{ozturkler2023smrd} & SURE-tuned DC weight and \textbf{early stopping against noise overfitting} & scalar & yes & yes\\
Lee \& Choi \cite{lee2026msc} & Plug-in measured-subspace correction wrapping six samplers & range/null & --- & yes\\
HFS-SDE \cite{cao2024hfssde}, CDiffMR \cite{huang2023cdiffmr} & Frequency-split diffusion processes (\emph{require training}) & k-space & --- & yes\\
Fast Diffusion EM \cite{laroche2024fastdiffem}, NA-NHMC \cite{xia2026nanhmc} & EM / marginalisation of blur kernel or scalar noise level & --- & yes & no\\
\midrule
\textbf{\fakgd{} (proposed framing)} & Per-frequency posterior variance in k-space, initialised from ACS, refined from residuals; frequency-gated DC & \textbf{k-space} & \textbf{yes} & \textbf{yes}\\
\bottomrule
\end{tabularx}
\end{center}

\paragraph{What remains defensible.}
\begin{enumerate}[label=(\roman*)]
  \item \emph{Basis matched to the forward operator.} For Cartesian MRI the measurement
        operator is diagonal in k-space up to coil sensitivities. A k-space-diagonal
        posterior variance therefore makes the Kalman update exact and costs $O(n)$.
        Peng et al.\ use a wavelet basis on natural images \cite{peng2024optimal}.
  \item \emph{Scan-specific, training-free estimation} from the ACS region and the
        residuals, with no calibration dataset.
  \item \emph{The empirical finding} that isotropic-DC samplers degrade as $T$ grows on
        multi-coil MRI. The literature search found no prior documentation for \pigdm-type
        samplers on MRI. SMRD \cite{ozturkler2023smrd} documents noise overfitting and
        early stopping, which is related evidence to cite. Burns \& Fridovich-Keil
        \cite{burns2026finite} give a theoretical lens: miscalibrated intermediate
        posterior spread makes samplers sensitive to stopping time.
\end{enumerate}

\paragraph{What must be conceded.} Frequency-progressive DC is conceptually close to
noise-conditioned frequency exposure \cite{tian2026frequency}. FPDC must be presented as
the MRI/k-space instantiation and compared against that paper's schedule, or derived
from the posterior-variance view (\S\ref{sec:theory}) rather than claimed as novel.
Non-isotropic posterior covariance in general is not new \cite{peng2024optimal,boys2024tmpd,rozet2024mmps}.

% ------------------------------------------------------------
\section{Proposed framing and a theory check}
\label{sec:theory}
% ------------------------------------------------------------
\paragraph{Equivalence.} \fakgd's gain
$K_i(t)=\sigma_t^2/(\sigma_t^2+\hat\sigma_i^2)$ is algebraically identical to a Kalman
gain $K_i=P_i/(P_i+\sigma_y^2)$ with \emph{white} measurement noise $\sigma_y^2$ and
per-frequency prior variance $P_i(t)=\sigma_y^2\sigma_t^2/\hat\sigma_i^2$. A radially
increasing ``noise'' profile is therefore equivalent to a radially \emph{decreasing}
posterior variance. That is physically expected, because image power falls off with
spatial frequency. The method survives the correction of its premise; only the
explanation changes. This also explains why the ACS coil-difference estimator (draft
Eq.~2) produces a radial profile on real data: coil differences contain sensitivity-weighted
\emph{signal}, which follows the power spectrum.

\paragraph{Linear-Gaussian reference model.} Assume a stationary Gaussian prior with
power spectrum $S_i$, observed through $x_t=x_0+\sigma_t\varepsilon$. Then
$\operatorname{Var}[x_{0,i}\mid x_t]=P_i(t)=S_i\sigma_t^2/(S_i+\sigma_t^2)$, and the
exact gain for a sampled frequency is
\begin{equation}
  K_i^{\star}(t)=\frac{P_i(t)}{P_i(t)+\sigma_y^2}
  =\Bigl(1+\tfrac{\sigma_y^2}{\sigma_t^2}+\tfrac{\sigma_y^2}{S_i}\Bigr)^{-1}.
  \label{eq:wiener}
\end{equation}
Compared with the isotropic gain $(1+\sigma_y^2/\sigma_t^2)^{-1}$, the only new term is
the spectral SNR $\sigma_y^2/S_i$. Two consequences follow.
\begin{enumerate}
  \item A \emph{principled, hyperparameter-free variant} exists: estimate $S_i$ from the
        ACS (radial power-law fit) and $\sigma_y^2$ from prewhitened background or
        repetitions. Call it \emph{Spectral-SNR DC}. It must be in the ablation.
  \item Eq.~\eqref{eq:wiener} does \emph{not} predict a progressive opening of high
        frequencies over time. FPDC's benefit therefore cannot be justified by the
        Gaussian model alone. It has to be justified empirically (C4 mechanism plots) or
        through approximation error of the Tweedie mean at large $\sigma_t$
        \cite{burns2026finite}. This is an honest, testable distinction. If Spectral-SNR DC
        recovers most of the gain, the paper becomes cleaner and more principled; if FPDC
        adds on top, that is evidence for an approximation-error mechanism.
\end{enumerate}
\emph{Before using this in the paper,} re-derive it with the EDM real/complex channel
scaling used in \code{multicoil.py} (the factor 2 in $K_{c,i}$) and with coil
sensitivities ($K$ becomes a small per-pixel system in SENSE).

% ------------------------------------------------------------
\section{Ranked changes}
\label{sec:changes}
% ------------------------------------------------------------
Each change lists the problem, the change, the literature grounding, the effort on one
A100, and a success criterion. Effort is in person-days (pd) plus GPU time.

% ---- Tier 1 -------------------------------------------------
\medskip\noindent\tier{Tier 1 --- blocking: decides whether there is a paper}

\change{C1}{Replace the noise premise; evaluate on real noise only}
\field{Problem.} Headline results depend on synthetic radial noise motivated by an
incorrect physical claim (\S\ref{sec:audit}). An MR-physics reviewer stops reading there.
\field{Change.}
(a) Remove the ``1--2 orders of magnitude noise'' claim. Adopt the posterior-variance
interpretation (\S\ref{sec:theory}).
(b) Run every main experiment on \emph{unmodified} raw k-space. Estimate the coil noise
covariance $\Psi$ from signal-free regions. The raw fastMRI readout is 2$\times$
oversampled, and the region outside the central 320 columns is mostly air. Prewhiten
\cite{kellman2005snr,pruessmann1999sense}.
(c) Keep synthetic noise only as a controlled \emph{SNR sweep}: add \emph{white,
$\Psi$-correlated} noise at chosen SNR levels, not radial noise.
(d) Add Spectral-SNR DC (Eq.~\ref{eq:wiener}) as a variant.
\field{Grounding.} Roemer \cite{roemer1990phased}; SENSE \cite{pruessmann1999sense};
SNR-unit reconstruction and noise-equivalent bandwidth \cite{kellman2005snr}; magnitude
statistics \cite{gudbjartsson1995rician,constantinides1997snr,ajafernandez2016noise}.
To our knowledge no major public raw dataset ships noise prescans: fastMRI, M4Raw,
CMRxRecon and Calgary-Campinas do not document them. Estimating $\Psi$ from background
is therefore standard practice and should be stated as a limitation.
\field{Effort.} 3--4 pd; under 1 GPU-day.
\field{Success (Gate G1).} On real-noise brain data at $R\in\{4,8\}$ and $T\in\{20,50,100\}$,
the frequency-aware variant beats the \emph{best tuned} isotropic baseline (C2). The
volume-level mean $\Delta$PSNR must be $\ge0.2$\,dB with a bootstrap 95\% CI excluding 0
for $T\ge50$.

\change{C2}{Correct and fairly tune the baselines}
\field{Problem.} The baseline is not \pigdm{}, and it shares \fakgd's hyperparameters
without its own search. The draft itself flags the ``tuning asymmetry'' (l.900--902).
Its degradation with $T$ may be a tuning artefact, and that would remove the central finding.
\field{Change.}
(a) Rename the current baseline \emph{Isotropic Kalman DC (IK-DC)}. It is the exact
ablation parent of \fakgd{} and should stay.
(b) Implement \pigdm{} as published (VJP guidance, $r_t$ schedule) \cite{song2023pigdm}.
(c) For every method, tune its own hyperparameters ($\sigma_y$, step weight, $\zeta$,
CG iterations) per $(R,T)$ on a disjoint calibration split, with an equal search budget.
InverseBench tunes every algorithm per task with a fixed protocol
\cite{zheng2025inversebench}; follow and cite it.
\field{Effort.} 4--6 pd; about 2 GPU-days for the sweeps.
\field{Success.} The IK-DC and \pigdm{} degradation with $T$ persists after tuning. If it
does not, C4 changes from ``samplers degrade'' to ``frequency-aware DC reaches the tuned
optimum without per-$T$ tuning''. That is weaker but still publishable.

\change{C3}{Diagnose the inert M-step; fix it or remove it}
\field{Problem.} Component 3 probably does nothing (\S\ref{sec:audit}), so the claimed
superadditive synergy (draft \S4.4, ablation table) cannot be attributed to it.
\field{Change.} On 32 real slices, compare \code{m\_step\_mode} $\in$ \{off, clamp, full\}
and log $\|\hat\sigma^2_t-\hat\sigma^2_{\mathrm{init}}\|/\|\hat\sigma^2_{\mathrm{init}}\|$
per step. If the change is $<1\%$, drop the component and the $\gamma$-family section.
If you keep it, re-specify it as an online estimate of the \emph{posterior variance}
$P_i(t)$ from residuals. Under \S\ref{sec:theory}, the expected squared residual is
$P_i+\sigma_y^2$, which is a principled moment estimator. Use a schedule that is not
switched off by $\alpha_t\to1$.
\field{Grounding.} Online moment estimation of Tweedie covariance \cite{hamidi2024cadps,boys2024tmpd};
EM-style adaptation \cite{laroche2024fastdiffem}.
\field{Effort.} 1--2 pd.
\field{Success.} Either a documented removal (simpler method, honest ablation) or a
measurable contribution of at least 0.05\,dB with a CI excluding 0.

% ---- Tier 2 -------------------------------------------------
\medskip\noindent\tier{Tier 2 --- core contribution}

\change{C4}{Make step-count degradation the central, mechanistically measured finding}
\field{Problem.} ``Scaling law'' is an overclaim, and the mechanism (``DC shock'') is
asserted rather than measured (draft l.448--451 admits this).
\field{Change.} Plot per-band k-space error of $\hat x_0(x_t)$ as a function of the
reverse step, for each sampler, with and without the frequency-aware module. Use 5
radial bands and median plus IQR over volumes. Also record the ``realised residual vs.\
predicted variance'' ratio per band, which tests the calibration of the implied
$P_i(t)$. Replace ``scaling law'' with ``step-robustness''. Extend the $T$ grid to
$\{10,20,50,100,200\}$ and report NFE explicitly.
\field{Grounding.} Noise overfitting and early stopping in diffusion MRI \cite{ozturkler2023smrd};
finite-sample failure modes \cite{burns2026finite}; few-step samplers showing that more
steps are not required \cite{pandey2024cpigdm}.
\field{Effort.} 3 pd; about 1 GPU-day.
\field{Success (Gate G2).} The per-band error curves show high-band error accumulating
in the isotropic samplers and not in the frequency-aware variant, on both anatomies.

\change{C5}{Multi-coil SENSE as the primary operator}
\field{Problem.} A coil-combined single-channel operator is non-standard for raw
fastMRI, and reviewers know this. It also makes the supervised gap uninterpretable.
\field{Change.} Use ESPIRiT maps \cite{uecker2014espirit} from the ACS, after prewhitening.
Extend the existing \code{run\_pigdm\_mc} and \code{run\_fakgd\_mc} paths
(\code{src/samplers/multicoil.py}) to all baselines. Report E2E-VarNet
\cite{sriram2020varnet} on exactly the same slices and masks as a supervised reference.
\field{Grounding.} Score-MRI parallel imaging \cite{chung2022scoremri}; DDS on multi-coil
MRI \cite{chung2024dds}; InverseBench's multi-coil fastMRI task with real raw measurements
\cite{zheng2025inversebench}; the prior from \cite{aali2025ambient}.
\field{Effort.} 4--6 pd; 2--3 GPU-days for the main grid.

\change{C6}{Plug-in evaluation across samplers, explicitly positioned}
\field{Problem.} A single-sampler comparison invites ``tuning asymmetry'' objections.
Recent work already proposes plug-in DC operators for MRI \cite{lee2026msc} and
frequency-conditioned exposure \cite{tian2026frequency}.
\field{Change.} Implement the frequency-aware gain as a drop-in for the DC step of
IK-DC, \pigdm{} \cite{song2023pigdm}, DDS \cite{chung2024dds}, DiffPIR \cite{zhu2023diffpir},
and, if time allows, DAPS \cite{zhang2025daps}. DPS \cite{chung2023dps} has no explicit DC
step, so report it as an unmodified reference. Main table: sampler $\times$
\{off, on\} $\times$ $T$. Add head-to-heads with (i) Tian et al.'s frequency-exposure
schedule ported to k-space and (ii) the measured-subspace operator \cite{lee2026msc}, if
its code is available. Otherwise reimplement it and label it as such.
\field{Effort.} 7--10 pd; 3--4 GPU-days.
\field{Success.} The module is positive for at least 3 of 4 samplers at $T\ge50$, with
volume-level CIs excluding 0, and never worse than $-0.05$\,dB.

\change{C7}{Leak-free protocol, volume-level statistics, compute reporting}
\field{Change.}
(a) Split by volume: calibration (tuning) / test. Use the official test split, or a
frozen held-out subset of validation volumes, for the test set.
(b) Freeze manifests (SHA-256 of the HDF5 files) before running the test.
(c) Primary statistic: paired volume-level mean $\Delta$ with bootstrap 95\% CIs; Wilcoxon
signed-rank across volumes; Holm correction across the main cells. Slice curves are
descriptive only.
(d) Report NFE, wall-clock time per slice on the A100, and peak memory.
(e) Metrics: PSNR and SSIM on RSS magnitude within a foreground mask, plus NMSE and HFEN.
Note Rician/noncentral-$\chi$ bias at low SNR \cite{gudbjartsson1995rician,constantinides1997snr}.
\field{Effort.} 2 pd. This is mostly reuse of the RaoSense provenance tooling.

% ---- Tier 3 -------------------------------------------------
\medskip\noindent\tier{Tier 3 --- clinical relevance and external validity (MIDL fit)}

\change{C8}{Anatomy-matched knee prior}
\field{Problem.} The knee results use a brain prior, and the cross-domain caveat runs
through the whole draft.
\field{Change.} First check whether the Ambient-DPS repository \cite{aali2025ambient} or
InverseBench \cite{zheng2025inversebench}, which uses an EDM prior on multi-coil fastMRI
knee, releases a knee checkpoint. Otherwise train an EDM \cite{karras2022edm} on fastMRI
knee training volumes. Our estimate for that is 3--5 A100-days at 320$\times$320; verify
with a 2-hour throughput test. Keep brain$\to$knee as a \emph{robustness} row.
\field{Effort.} 1--2 pd plus 0--5 GPU-days. Start in week 1 so training runs in parallel.

\change{C9}{Real low-SNR evidence}
\field{Problem.} A ``noise-aware'' method must be shown where noise matters.
\field{Change.}
\begin{itemize}
  \item \textbf{M4Raw} \cite{lyu2023m4raw}: 0.3\,T, 4-channel brain, 183 subjects,
        T1w/T2w $\times3$ and FLAIR $\times2$ repetitions, CC~BY~4.0. Differences of
        repetitions give an empirical noise estimate, and averages give a higher-SNR
        reference. Use the brain prior as-is (robustness) and, if time allows, fine-tune
        it.
  \item \textbf{fastMRI Prostate} \cite{tibrewala2024prostate}: T2w is stored as 3
        separate averages; averages 1 and 3 sample the same (odd) lines. This gives a
        second real-noise estimate at 3\,T with many channels. Optional, as a third
        anatomy.
  \item \textbf{SNR sweep} on fastMRI brain using $\Psi$-correlated white noise
        (C1c), reported in SNR units \cite{kellman2005snr}.
\end{itemize}
\field{Effort.} 4--5 pd; 1--2 GPU-days.
\field{Success (Gate G3).} The gain over the tuned baseline increases as SNR decreases.

\change{C10}{Pathology-region and hallucination analysis}
\field{Change.} Use fastMRI+ boxes \cite{zhao2022fastmriplus} (brain: 7{,}570 boxes, 30
categories; knee: 16{,}154 boxes, 22 categories). Report ROI-restricted PSNR/SSIM/NMSE
against size-matched normal patches, and hallucination maps \cite{bhadra2021hallucination}.
Show a frozen, quantile-chosen qualitative panel, not hand-picked slices. Add deep
feature distances, which track radiologist scores better than PSNR/SSIM
\cite{adamson2023dfd}. Motivate this with the documented hallucinations in the 2020
fastMRI challenge \cite{muckley2021fastmri}.
\field{Effort.} 3--4 pd; under 1 GPU-day (reuses reconstructions). This is the main lever
for the ``clinical relevance'' requirement.

\change{C11}{Corrupted-sample robustness [stretch; requires C3 to work]}
\field{Idea.} Per-frequency trust that adapts from residuals should down-weight corrupted
samples: RF spikes / herringbone \cite{campbellwashburn2016spike,jin2017aloha} and
motion-corrupted phase-encode lines \cite{levac2024motion}.
\field{Change.} Simulate on \emph{multi-coil raw} k-space with our own implementation.
TorchIO's RandomSpike, RandomMotion and RandomGhosting \cite{perezgarcia2021torchio} operate on
single-channel k-space synthesised from magnitude images, so use them only as reference
implementations. EMI at low field \cite{srinivas2022editer,liu2021ulf} is out of scope
(no public raw EMI data).
\field{Effort.} 4--6 pd. Cut if the timeline slips.

\medskip\noindent\tier{Tier 4 --- valuable, but realistically after 4 December}

\change{C12}{Downstream task evaluation}
SKM-TEA \cite{desai2021skmtea} offers cartilage segmentation and T2 maps from GE 3D qDESS,
and K2S \cite{tolpadi2023k2s} offers segmentation from undersampled knee data. Caveat:
both ``fully sampled'' references are ARC-filled. K2S reports no correlation between
reconstruction and segmentation metrics, which is a strong motivation for a journal
extension.

\change{C13}{Uncertainty calibration}
With several posterior samples per slice, report the correlation between per-pixel
spread and error, and per-band coverage. This follows naturally from the
posterior-variance framing but needs about 8$\times$ more compute.

\medskip\noindent\tier{Cross-cutting}

\change{C14}{Rewrite for MIDL}
\begin{itemize}
  \item Use the official MIDL template, a 10-page main text, and author names (single-blind).
  \item Remove the CS231N-specific sections (Contributions, ``Dataset and Features'',
        GenAI statement in its current form).
  \item Put a clinical-relevance paragraph in the introduction (scan time; artefact
        risk; hallucination concerns \cite{muckley2021fastmri}).
  \item Fix every finding in \S\ref{sec:audit}: ``scaling law''$\to$``step-robustness'';
        stochastic re-noising, not DDIM; one oracle specification or none.
  \item Write the related-work section from \S\ref{sec:lit}, positioning explicitly
        against \cite{peng2024optimal,tian2026frequency,ozturkler2023smrd,lee2026msc}.
  \item \textbf{LLM policy:} the authors write all substantive text, with LLM use limited
        to language polishing \cite{midl2027cfp}. Keep a change log of who wrote what.
  \item Release code and frozen manifests (encouraged by MIDL).
\end{itemize}

% ------------------------------------------------------------
\section{Datasets and noise models: selection}
\label{sec:data}
% ------------------------------------------------------------
\begin{center}
\small
\begin{tabularx}{\linewidth}{@{}l L c@{}}
\toprule
Dataset & Facts relevant to us & Use\\
\midrule
fastMRI brain/knee \cite{zbontar2018fastmri,knoll2020fastmri} & Raw multi-coil; Siemens 1.5/3\,T; knee 15 channels; readout 2$\times$ oversampled; no documented noise prescans & \textbf{core}\\
fastMRI+ \cite{zhao2022fastmriplus} & Pathology boxes for fastMRI brain/knee & \textbf{core} (C10)\\
M4Raw \cite{lyu2023m4raw} & 0.3\,T, 4-channel brain, 183 subjects, repetitions, CC~BY~4.0 & \textbf{core} (C9)\\
fastMRI Prostate \cite{tibrewala2024prostate} & 3\,T Siemens Vida; 3 separate T2w averages; 10--30 channels & optional (C9)\\
Calgary-Campinas MC \cite{beauferris2022calgary} & GE 3\,T, 12/32-channel 3D T1 brain; vendor shift & optional\\
CMRxRecon \cite{wang2024cmrxrecon} & Siemens 3\,T, 32-channel cardiac, CC~BY~4.0 & journal\\
SKM-TEA \cite{desai2021skmtea}, K2S \cite{tolpadi2023k2s} & GE 3\,T knee; segmentation, T2; references ARC-filled & journal (C12)\\
mridata.org \cite{ong2018mridata} & Stanford 3D FSE knees (20 subjects, 8 channels) & sanity check only\\
\bottomrule
\end{tabularx}
\end{center}

\paragraph{Noise models we will (and will not) use.}
\begin{itemize}
  \item \textbf{Use:} thermal noise, white per channel and coil-correlated through $\Psi$,
        with prewhitening \cite{roemer1990phased,kellman2005snr}; controlled SNR reduction
        by adding $\Psi$-correlated white noise; real low-field noise (M4Raw); magnitude
        statistics for evaluation \cite{gudbjartsson1995rician,constantinides1997snr}.
  \item \textbf{Stretch:} sparse k-space spikes; motion-corrupted lines (C11).
  \item \textbf{Do not use as motivation:} ``frequency-dependent thermal noise''. Readout
        anti-aliasing roll-off affects image-domain position and the noise-equivalent
        bandwidth \cite{kellman2005snr}, not per-sample k-space variance across spatial
        frequency.
\end{itemize}

% ------------------------------------------------------------
\section{Timeline to 4 December 2026, with gates}
\label{sec:timeline}
% ------------------------------------------------------------
\begin{center}
\small
\begin{tabularx}{\linewidth}{@{}l L l@{}}
\toprule
Week (2026) & Work & Gate\\
\midrule
1: Oct 5--11 & C3 diagnosis; real-noise brain run (C1); IK-DC vs.\ tuned IK-DC (C2a/c); start knee-prior throughput test or training (C8) & \textbf{G1 (Oct 12)}\\
2: Oct 12--18 & Implement \pigdm{} and multi-coil baselines (C2b, C5); freeze splits and manifests (C7) & \\
3: Oct 19--25 & DDS, DiffPIR, DAPS plug-ins (C6); mechanism instrumentation (C4) & \\
4: Oct 26--Nov 1 & Main grid on brain; mechanism plots & \textbf{G2 (Nov 1)}\\
5: Nov 2--8 & Knee grid (matched prior); M4Raw and SNR sweep (C9) & \textbf{G3 (Nov 8)}\\
6: Nov 9--15 & fastMRI+ ROI and hallucination analysis (C10); C11 only if ahead & \\
7: Nov 16--22 & Ablations: Spectral-SNR DC vs.\ FPDC vs.\ both; M-step; figures frozen & \\
8: Nov 23--30 & Writing by the authors (C14); internal review; \textbf{abstract registration Nov 30} & \\
9: Dec 1--4 & Final pass; submit by \textbf{Dec 4 AoE} & \\
\bottomrule
\end{tabularx}
\end{center}

\paragraph{Compute budget (one A100).} The inference grid is roughly 200 test slices
$\times$ 2 anatomies $\times$ 5 samplers $\times$ \{off, on\} $\times$ 5 $T$-values
$\times$ 2 $R$, plus M4Raw and the SNR sweep. At a few seconds per slice for $T\le100$,
this totals 3--6 GPU-days. Knee-prior training (0--5 days) is the only large item and
runs concurrently. Measure per-slice runtime in week 1 and re-plan if it is above
10\,s.

% ------------------------------------------------------------
\section{Risks and fallbacks}
% ------------------------------------------------------------
\begin{center}
\small
\begin{tabularx}{\linewidth}{@{}L L@{}}
\toprule
If\ldots & Then\ldots\\
\midrule
G1 fails: no gain over a tuned isotropic baseline on real noise & Do not submit a methods paper to the main track. Write a negative/analysis note (``frequency-aware DC under real noise'') for the MIDL short-paper track (spring 2027), and keep FA-KGD as a course report.\\
Baseline degradation with $T$ disappears after tuning & Reframe C4: the frequency-aware gain removes the need for per-$T$ tuning. Report tuning cost as the endpoint.\\
Spectral-SNR DC explains all of the gain & Lead with the principled variant; FPDC becomes an ablation. This is a cleaner paper.\\
M-step inert (C3) & Remove it and the $\gamma$-family section; drop C11.\\
Knee prior not ready by week 5 & Report brain (matched) and M4Raw; keep knee as a cross-domain robustness row with an explicit caveat.\\
Tian et al.\ \cite{tian2026frequency} schedule matches FPDC on MRI & Concede FPDC; the contribution becomes the k-space posterior-variance gain plus the MRI step-robustness analysis.\\
\bottomrule
\end{tabularx}
\end{center}

% ------------------------------------------------------------
\section{Target paper outline (10 pages, MIDL main track)}
% ------------------------------------------------------------
\begin{enumerate}
  \item \textbf{Introduction} (1\,p): clinical motivation; failure of isotropic DC as
        $T$ grows; contributions (k-space posterior-variance gain; step-robustness
        analysis; plug-in evaluation on real multi-coil and low-field data).
  \item \textbf{Related work} (0.75\,p): \S\ref{sec:lit}.
  \item \textbf{Method} (2\,p): linear-Gaussian derivation (Eq.~\ref{eq:wiener}); ACS
        spectral estimate; FPDC as an approximation-error correction; residual-based
        refinement (if C3 survives); plug-in form for each sampler.
  \item \textbf{Experimental setup} (1.25\,p): data, prewhitening, splits, tuning
        protocol, metrics, statistics, compute.
  \item \textbf{Results} (3.5\,p): main plug-in table; step-robustness and mechanism
        figure; low-SNR and M4Raw; pathology ROI and hallucinations; ablations.
  \item \textbf{Discussion and limitations} (1\,p): $\Psi$ from background (no
        prescans); retrospective undersampling; no reader study; brain-trained priors
        where applicable.
  \item \textbf{Appendix:} derivations, per-volume results, full tuning grids,
        qualitative panels.
\end{enumerate}

\paragraph{Items to verify before citing.} Before any reference marked UNVERIFIED in the
bibliography is cited, check its author lists and venue on the primary source. Also
check the InverseBench MRI numbers against its tables. Read Blumenthal et al.\
\cite{blumenthal2025pula} to see whether its preconditioner is k-space diagonal, because
that would be a further novelty risk.

\bibliographystyle{plainnat}
\bibliography{references}

\end{document}
